\begin{textsample}{POS Dim 4 – 2010\_c4 – Score -1.00 – 2010\_c4/enade\_2010\_tec\_rad\_711bf0db15.txt}  \label{ex:f4_pos_295}
CAPÍTULO 4 PERCEPÇÃO DA PROVA As análises feitas neste capítulo tratam das percepções dos estudantes, separados por grupos de Ingressantes e Concluintes, sobre a prova aplicada no ENADE/2010 na Área de Tecnologia em Radiologia.
Estas percepções foram mensuradas por meio de nove questões que avaliaram desde o grau de dificuldade da prova até o tempo gasto para concluí -la.
As percepções sobre a prova foram relacionadas com o desempenho dos alunos e com a Grande Região de funcionamento do curso.
O desempenho dos alunos foi classificado em dois níveis : quarto inferior e quarto superior.
Para tanto, esse desempenho foi colocado em ordem ascendente.
O percentil 25, P25, também conhecido como primeiro quartil, é a nota de desempenho que deixa um quarto ( 25,0 \% ) dos valores observados abaixo e três quartos acima.
A Figura 1 apresenta uma ilustração destes conceitos.
O quarto inferior de desempenho é composto pelas notas abaixo do primeiro quartil.
Já o percentil 75, P75, também conhecido como terceiro quartil, é o valor para o qual há três quartos ( 75,0 \% ) dos dados abaixo e um quarto acima dele.
O quarto superior de desempenho é composto pelas notas iguais ou acima do terceiro quartil.
O percentil 50, P50, também conhecido como mediana, é o valor que divide as notas em dois conjuntos de igual tamanho.
Vale ressaltar que percentis, quartis e medianas são pontos, e não obrigatoriamente pertencentes ao conjunto original de dados.
Ao passo que os quartos são subconjuntos dos dados originais.
A seguir, serão apresentados gráficos com resultados selecionados, relativos às nove questões avaliadas por grupos de estudantes.
As Tabelas no Anexo II apresentam os valores absolutos e a distribuição percentual das alternativas válidas das nove questões, segundo o mesmo recorte de desempenho dos alunos e Grande Região de funcionamento do curso. 4.1 GRAU DE DIFICULDADE DA PROVA 4.1.1 Componente de Formação Geral Na questão “ Qual o grau de dificuldade desta prova na parte de Formação Geral? ”, 28,8 \% do grupo de Ingressantes e 27,8 \% do grupo de Concluintes optaram pelas alternativas difícil ou muito difícil.
Entretanto, para mais da metade dos estudantes, tanto Ingressantes ( 60,8 \% ) quanto Concluintes ( 62,4 \% ), o Componente de Formação Geral da prova foi considerado com grau de dificuldade médio ( Gráfico 4.1 e Tabelas II.1 e II.10 no Anexo II ).
O percentual de Ingressantes que consideraram a prova difícil ou muito difícil foi mais elevado no quarto inferior ( 36,9 \% ) do que no superior ( 21,4 \% ), o que também foi verificado no grupo dos Concluintes, com percentual de participantes maior no quarto inferior ( 33,3 \% ), do que no quarto superior ( 22,4 \% ).
Semelhantemente ao ocorrido para os Ingressantes como um todo, para os quartos de menor e maior desempenho deste grupo, a alternativa modal para esta pergunta foi médio, com respectivamente 54,9 \% e 63,4 \% dos respondentes.
Já entre os Concluintes, estas percentagens foram de, respectivamente, 61,0 \% e 62,7 \%.
Considerando -se as alternativas escolhidas pelos estudantes de cada Grande Região, a Centro-Oeste foi aquela em que houve maior proporção de Ingressantes que avaliaram o Componente de Formação Geral da prova como difícil ou muito difícil : 32,1 \% tiveram essa percepção.
Os Concluintes emitiram essa opinião, em maior medida, também na região Centro-Oeste ( 31,4 \% ).
Nas Grandes Regiões, a proporção de Ingressantes que consideraram o Componente de Formação Geral da prova como sendo de grau de dificuldade médio esteve localizada entre 59,3 \% na região Sudeste e 66,5 \% na região Norte.
Optando pela alternativa de grau médio, por Grande Região, no grupo dos Concluintes os percentuais estiveram entre 60,0 \% na região Centro-Oeste e 68,1 \% na região Sul. 4.1.2 Componente de Conhecimento Específico Ao responderem à questão “ Qual o grau de dificuldade desta prova na parte de Componente Específico? ”, 56,0 \% do grupo de Ingressantes e 38,2 \% do grupo de Concluintes classificaram -na como difícil ou muito difícil.
O Componente de Conhecimento Específico da prova foi considerado com grau de dificuldade médio por 40,4 \% dos Ingressantes e por 57,0 \% dos Concluintes ( Gráfico 4.2 e Tabelas II.2 e II.11 no Anexo II ).
Considerando -se a percepção sobre o grau de dificuldade do Componente de Conhecimento Específico da prova de acordo com o desempenho dos estudantes, observa -se que 58,4 \% dos Ingressantes que obtiveram desempenho situado no quarto superior classificaram -na como difícil ou muito difícil.
Dos Ingressantes que obtiveram desempenho situado no quarto inferior, 50,7 \% optaram pela mesma alternativa.
Para os Concluintes, o percentual de estudantes que consideraram o Componente de Conhecimento Específico da prova difícil ou muito difícil foi um pouco maior entre os que obtiveram desempenho no quarto inferior ( 36,2 \% ) do que os que obtiveram desempenho no quarto superior ( 35,9 \% ).
Para os Ingressantes do quarto de melhor desempenho, a alternativa modal para esta pergunta foi médio, com 63,4 \%.
Já entre os Concluintes, a alternativa modal para a Questão 2 foi médio para ambos os quartos, com 60,1 \% do quarto inferior e 59,1 \% do superior.
Entre as Grandes Regiões brasileiras, a proporção de Ingressantes que avaliaram o Componente de Conhecimento Específico da prova como difícil ou muito difícil variou de 42,5 \% ( Norte ) a 66,4 \% ( Sul ).
O grupo de Concluintes avaliou esse componente como difícil ou muito difícil com menor frequência, com o mínimo alcançado na região Norte, 27,4 \%, e o máximo na região Centro-Oeste, com 44,6 \% dos respondentes.
O percentual de Ingressantes que classificaram o grau de dificuldade como médio, no Componente de Conhecimento Específico da prova variou de 31,3 \% ( Sul ) a 52,5 \% ( Norte ).
Dentre os Concluintes, a percepção do mesmo grau de dificuldade variou de 49,0 \% ( Centro-Oeste ) a 71,0 \% ( Norte ). 4.2 EXTENSÃO DA PROVA EM RELAÇÃO AO TEMPO TOTAL Indagados quanto à extensão da prova, em relação ao tempo total oferecido para a sua resolução, os estudantes apontaram, com maior incidência, a alternativa que considerava a extensão adequada, para ambos os grupos e todas as agregações consideradas ( Gráfico 4.3 e Tabelas II.3 e II.12 no Anexo II ).
No grupo de Ingressantes, a maioria, 71,9 \%, respondeu ser a extensão da prova adequada, ao passo que, no grupo de Concluintes, este valor foi menor : 66,2 \%.
Para 24,9 \% dos Ingressantes a extensão do exame foi longa ou muito longa.
Essa mesma consideração foi feita por 29,3 \% dos Concluintes.
No grupo de Ingressantes, a maioria, tanto para o quarto de desempenho inferior ( 68,6 \% ) quanto para o de melhor desempenho ( 75,8 \% ) respondeu ser a extensão da prova adequada.
Já entre os Concluintes, esses valores foram menores, situados em cada quarto em, respectivamente, 59,0 \% e 69,8 \%.
Considerando -se ainda os quartos extremos de desempenho, observa -se que a proporção de estudantes que consideraram a prova longa ou muito longa, em relação ao tempo total destinado à sua resolução, acompanhou a percepção de Ingressantes e Concluintes.
Para o grupo de Ingressantes, entre os estudantes no quarto de desempenho inferior, 27,9 \% declararam ter esta percepção, ao passo que, entre aqueles no quarto de melhor desempenho, esta proporção foi 21,6 \%.
O mesmo pode ser constatado entre os Concluintes : 34,8 \% no quarto de menor desempenho e 25,6 \% no quarto de maior desempenho declararam ter tido esta percepção.
Na análise por Grandes Regiões, o percentual de Ingressantes que responderam ser a extensão da prova adequada variou de 70,8 \% no Sudeste a 74,4 \%, no Norte e Nordeste.
Já entre os Concluintes, estes valores estavam compreendidos entre 63,6 \% ( Centro-Oeste ) e 74,0 \% ( Sul ).
Aqueles Ingressantes que avaliaram a prova como longa ou muito longa em relação ao tempo total destinado à sua resolução variou de 21,2 \% na região Norte a 26,6 \% na região Sudeste.
Os percentuais, quanto ao mesmo aspecto, entre os Concluintes, foram de 22,2 \% na região Sul, indo até 32,1 \% na região Centro-Oeste. 4.3 COMPREENSÃO DOS ENUNCIADOS DAS QUESTÕES 4.3.1 Componente de Formação Geral Com relação aos enunciados das questões do Componente de Formação Geral, as opiniões foram positivas, com percentuais de 74,2 \% para Ingressantes e 74,7 \% para Concluintes que consideraram claros e objetivos os enunciados de todas ou da maioria das questões ( Gráfico 4.4 e Tabelas II.4 e II.13 no Anexo II ).
Ao se comparar as percepções dos estudantes segundo o desempenho, observou -se que, entre os Ingressantes do quarto superior, a clareza e objetividade de todos ou da maioria dos enunciados das questões foram percebidas por 84,2 \%.
A proporção apurada para os Concluintes desse mesmo quarto foi de 82,0 \%.
No quarto de desempenho inferior, tal avaliação foi emitida por 66,2 \% dos Ingressantes e por 67,3 \% dos Concluintes.
A partir dessas estatísticas, é possível afirmar que todos ou a maioria dos enunciados das questões do Componente de Formação Geral estavam claros e objetivos para a maior parte dos respondentes, segundo os quartos de desempenho considerados, embora os que se situaram no quarto superior tenham declarado esta percepção com maior incidência.
Na análise regional, observa -se que as opiniões de que todos ou a maioria dos enunciados das questões do Componente de Formação Geral estavam claros e objetivos variaram, entre os Ingressantes, de 71,9 \% no Sudeste a 81,1 \% no Sul.
Entre os Concluintes os valores correspondentes variaram de 72,0 \% no Sudeste a 80,8 \% no Norte. 4.3.2 Componente de Conhecimento Específico Com relação aos enunciados das questões do Componente de Conhecimento Específico da prova, para os Ingressantes da Área de Tecnologia em Radiologia, a clareza e a objetividade foram percebidas em todos ou na maioria destes por 69,0 \%.
Para o grupo dos Concluintes, foram percebidas por 75,6 \%.
A proporção de estudantes que considerou os enunciados das questões claros e objetivos foi mais elevada nos quartos superiores : 79,0 \% dos Ingressantes e 83,3 \% dos Concluintes, se comparados aos 62,9 \% e aos 67,2 \% dos mesmos grupos, respectivamente, que estiveram no quarto inferior de desempenho ( Gráfico 4.5 e Tabelas II.5 e II.14 no Anexo II ).
Entre as Grandes Regiões brasileiras, os percentuais de Concluintes que consideraram claros e objetivos todos ou a maioria dos enunciados das questões do Componente de Conhecimento Específico da prova foram mais elevados na região Sul ( 81,1 \% ), e menos expressivo na região Centro-Oeste ( 74,3 \% ).
Quanto aos Ingressantes, o percentual mais baixos ocorreu na região Sudeste onde 66,6 \% dos estudantes manifestaram -se da mesma forma.
Os índices mais altos desta percepção foram encontrados, para os Ingressantes, na região Sul ( 75,5 \% ). 4.4 SUFICIÊNCIA DAS INFORMAÇÕES/INSTRUÇÕES FORNECIDAS Ao avaliarem as informações/instruções fornecidas para a resolução das questões, 72,4 \% do grupo dos Ingressantes e 79,0 \% do grupo dos Concluintes respondentes de todo o Brasil afirmaram que estas eram até excessivas ou suficientes em todas ou na maioria das questões ( Gráfico 4.6 e Tabelas II.6 e II.15 no Anexo II ).
Levando -se em conta o desempenho dos participantes, notam -se diferenças entre as opiniões de estudantes, Ingressantes e Concluintes, dos quartos inferior e superior de desempenho, como mostra o Gráfico 4.6.
O percentual de participantes que avaliaram as informações/instruções como excessivas ou suficientes em todas ou na maioria das questões foi mais elevado nos quartos superiores, correspondendo a 81,7 \% dos Ingressantes e a 88,9 \% dos Concluintes, percentuais superiores à média nacional para cada um dos dois grupos.
Já nos quartos inferiores, a suficiência das informações/instruções declarada como excessiva em todas ou na maioria das questões foi percebida por 66,1 \% dos Ingressantes e por 72,6 \% dos Concluintes. cada um dos dois grupos.
Já nos quartos inferiores, a suficiência das informações/instruções declarada como excessiva em todas ou na maioria das questões foi percebida por 66,1 \% dos Ingressantes e por 72,6 \% dos Concluintes.
Quanto à distribuição de respondentes nas Grandes Regiões, observa -se que as proporções de estudantes que perceberam como excessivas ou suficientes em todas ou na maioria das questões as informações/instruções fornecidas para a resolução das mesmas superaram as médias nacionais, em maior medida, entre os Ingressantes da região Norte ( 77,2 \% ).
Entre os Concluintes, a superação da média nacional foi mais expressiva na região Sul ( 86,5 \% ). 4.5 DIFICULDADE ENCONTRADA AO RESPONDER À PROVA Perguntados sobre as dificuldades com as quais se depararam ao responder à prova, 61,5 \% do grupo de Ingressantes e 14,0 \% do grupo de Concluintes apontaram o desconhecimento do conteúdo.
Concluintes, por sua vez, apontaram com maior frequência forma diferente de abordagem do conteúdo como dificuldade, mencionada por 52,1 \%.
Esta alternativa foi mencionada por uma parcela menor dos Ingressantes, 24,9 \%.
A falta de motivação para fazer a prova foi a dificuldade para 16,5 \% dos Concluintes e para 5,2 \% dos Ingressantes.
Considerando -se os respondentes em todo o Brasil, 6,7 \% dos Ingressantes e 14,8 \% dos Concluintes afirmaram que não tiveram qualquer tipo de dificuldade para responder à prova ( Tabelas II.7 e II.16 no Anexo II ).
Com relação aos quartos de desempenho, o desconhecimento do conteúdo foi a dificuldade modal entre os Ingressantes, apontada por 65,8 \% destes no quarto superior e por 55,8 \% no inferior.
Entre os Concluintes, os percentuais de estudantes com a mesma percepção foi inferior ao encontrado entre os Ingressantes, 15,8 \% no quarto superior e 14,0 \% no quarto inferior.
A alternativa modal pelo grupo de Concluintes, quando desagregada pelo desempenho, foi que a dificuldade encontrada era a forma diferente de abordagem do conteúdo : 46,6 \% no quarto inferior e 51,6 \% do quarto superior assim o responderam.
Na análise por Grandes Regiões, o percentual de Ingressantes que apontaram o desconhecimento do conteúdo como dificuldade ao responder à prova variou de 48,7 \%, na região Norte a 66,2 \%, na Sul.
Entre os Concluintes, tal percentual foi mais baixo, variando de 8,0 \% na região Norte a 16,1 \%, na Nordeste.
A forma diferente de abordagem do conteúdo foi citada, entre os Concluintes, por percentuais que variaram entre 46,5 \% ( região Nordeste ) e 60,0 \% ( região Norte ).
Em contrapartida, os percentuais relacionados a essa opção entre os Ingressantes foi mais baixo, variando de 20,8 \% ( região Sul ) a 31,2 \% ( região Norte ).
O percentual de Ingressantes que citou a falta de motivação como dificuldade superou a média nacional, em maior medida, na região Norte ( 7,0 \% ).
Entre os Concluintes esse percentual variou entre 10,4 \% ( região Norte ) e 19,6 \% ( região Nordeste ).
Houve percentual acima da média nacional dos Ingressantes nas regiões Nordeste ( 9,1 \% ) e Norte ( 11,1 \% ), com declarações de que não tiveram qualquer dificuldade para responder à prova.
Entre os Concluintes, esta situação aconteceu nas regiões Norte ( 16,0 \% ), Nordeste ( 16,1 \% ) e Sul ( 21,1 \% ).
O Gráfico 4.7 apresenta os percentuais de estudantes, Ingressantes e Concluintes, que apontaram o desconhecimento do conteúdo como dificuldade percebida ao responder à prova. 4.6 CONTEÚDOS DAS QUESTÕES OBJETIVAS DA PROVA Ao analisarem os conteúdos das questões objetivas da prova, a maior parte dos Ingressantes, 76,4 \% do total nacional, afirmou que não estudou ainda a maioria desses conteúdos, enquanto entre os Concluintes, 8,3 \% expressaram a mesma opinião.
Por outro lado, a maioria dos Concluintes ( 56,9 \% ) afirmou ter estudado e aprendido muitos ou todos os conteúdos.
Para os Ingressantes, o percentual correspondente a essas alternativas foi 10,1 \% ( Gráfico 4.8 e Tabelas II.8 e II.17 no Anexo II ).
Considerando -se separadamente as opiniões de estudantes dos quartos de desempenho inferior e superior, observa -se que, para os Ingressantes, o percentual de participantes que afirmaram que não estudaram ainda a maioria desses conteúdos, foi mais elevado no quarto superior ( 77,3 \% ) do que no inferior ( 73,3 \% ).
Para os Concluintes do quarto inferior, 12,5 \% ofereceram essa resposta, sendo 4,7 \% os do quarto superior com esta percepção.
Tendo em conta os quartos superiores para os dois grupos, 72,6 \% dos Concluintes e 13,4 \% dos Ingressantes afirmaram ter estudado e aprendido muitos ou todos os conteúdos.
Na análise por Grande Região, a proporção de Ingressantes que afirmou que não estudou ainda a maioria desses conteúdos superou a média nacional nas regiões Nordeste ( 77,2 \% ), Centro-Oeste ( 80,9 \% ) e Sul ( 81,8 \% ).
Já entre os Concluintes, as regiões Nordeste ( 11,3 \% ) e Centro-Oeste ( 9,5 \% ) apresentaram percentuais maiores do que a média nacional para esta alternativa.
Em todas as regiões, Concluintes afirmaram ter estudado e aprendido muitos ou todos os conteúdos, com percentuais variando entre 54,1 \% na região Norte e 70,0 \% na Sul. 4.7 TEMPO GASTO PARA CONCLUIR A PROVA Ao responderem sobre o tempo de conclusão da prova, mais da metade dos estudantes afirmou ter gasto entre duas e quatro horas.
Considerando -se separadamente as percepções do grupo de Ingressantes e do grupo de Concluintes, os percentuais, em termos nacionais, foram de 50,7 \% a 62,4 \%, respectivamente ( Gráfico 4.9 e Tabelas II.9 e II.18 no Anexo II ).
Uma vez sendo analisadas as alternativas oferecidas pelos estudantes que se situaram nos quartos superior e inferior, observa -se que uma maior proporção de Ingressantes no quarto superior declarou ter gasto entre duas e quatro horas para concluir a prova quando comparados com os do quarto inferior, respectivamente 65,5 \% e 34,9 \%.
A mesma situação foi verificada no grupo dos Concluintes, no qual 74,3 \% daqueles no quarto superior e 48,4 \% daqueles no quarto inferior declaram ter gasto entre duas e quatro horas para concluir a prova.
Nas cinco Grandes Regiões brasileiras, os Ingressantes que utilizaram entre duas e quatro horas para finalizar a prova superaram o percentual nacional nas regiões Centro-Oeste e Nordeste ( ambas com 52,5 \% ), e Norte ( 55,5 \% ).
Entre os Concluintes, o mesmo foi verificado entre os estudantes das regiões Sul ( 62,5 \% ), Centro-Oeste ( 64,2 \% ) e Norte ( 64,8 \% ).
A região Nordeste foi a que apresentou o menor percentual de Concluintes que dispensaram entre duas e quatro horas para concluir a prova : 61,2 \%, como mostra o Gráfico 4.9.

% matched lemmas: 
\end{textsample}
